\section{Where to perturb a ViT}
\label{sec:results}

We walk from the original placement to PSBD-TM on the same backdoored ViT-B/16 models, every number read at the rate the original rule assigns. \Cref{app:walk} gives each step in full, and the per-model numbers are in \cref{app:results}.

\paragraph{Dropout alone does not carry the result.}
PSBD-RD, dropout on the stream after both residual additions, reaches a mean AUROC of \StaircaseResidualPostResidualAuroc{} and reads below chance on \StaircaseResidualPostResidualBelowChance{} of \StaircaseResidualPostResidualN{} models. It is close to perfect on BadNets, Blend and LF at the higher poison rates and close to a coin on TaCT and on most models at \PanelRateLow{} poisoning. Moving the same dropout to the attention input, before its LayerNorm, reaches \StaircaseDropoutSitesBeforeAttentionNormAuroc{}, a paired gain of \StaircaseDropoutSitesBeforeAttentionNormGain{} [\StaircaseDropoutSitesBeforeAttentionNormGainLow{}, \StaircaseDropoutSitesBeforeAttentionNormGainHigh{}] whose interval touches 0. The site matters, but what is injected there matters more.

\subsection{Masking tokens instead of dropping coordinates}
\label{sec:results:operator}

\Cref{tab:staircase-operators} holds the site at the attention input and changes what is injected there. It then reads the same perturbations at the other sites, including Gaussian noise on both sides of the MLP's LayerNorm.

% generated by scripts/paper/tab_staircase.py from results/coverage/coverage.json, results/<folder>/psbd_metrics.json (54 models) at 2026-09-29T18:11:59.373806+00:00, commit d73d3277d007fd75b3a4166a7911b1e07a2761b7-dirty. Do not edit by hand.
\begin{table*}[htbp]
\centering
\small
\caption{Perturbation operators at the attention input of ViT-B/16, and token masking at the other sites, all 12 blocks.}
\label{tab:staircase-operators}
\begin{tabular}{lrrrrl}
\toprule
 & & & \multicolumn{2}{c}{TPR@FPR} & \\
\cmidrule(lr){4-5}
Placement & n & AUROC & 10\% & 20\% & Gain [95\% CI] \\
\midrule
token mask, attention input, before the norm & 54 & 0.963 & 0.887 & 0.916 & -- \\
channel mask, attention input, before the norm & 54 & 0.921 & 0.768 & 0.840 & $-$0.042 [$-$0.068, $-$0.018] \\
gaussian, attention input, before the norm & 54 & 0.836 & 0.595 & 0.705 & $-$0.127 [$-$0.191, $-$0.071] \\
dropout, attention input & 54 & 0.903 & 0.715 & 0.815 & $-$0.060 [$-$0.096, $-$0.031] \\
token mask, attention output before the add & 54 & 0.956 & 0.863 & 0.918 & $-$0.007 [$-$0.030, +0.016] \\
token mask, both sublayer inputs & 54 & 0.928 & 0.793 & 0.849 & $-$0.035 [$-$0.067, $-$0.009] \\
token mask, MLP input, before the norm & 54 & 0.838 & 0.556 & 0.731 & $-$0.125 [$-$0.173, $-$0.084] \\
gaussian, MLP input, before the norm & 35 & 0.697 & 0.397 & 0.557 & $-$0.286 [$-$0.397, $-$0.182] \\
gaussian, MLP input, after the norm & 54 & 0.904 & 0.781 & 0.833 & $-$0.059 [$-$0.106, $-$0.018] \\
token mask, after the attention add & 54 & 0.805 & 0.528 & 0.635 & $-$0.158 [$-$0.225, $-$0.099] \\
channel mask, MLP neurons & 54 & 0.833 & 0.597 & 0.717 & $-$0.130 [$-$0.179, $-$0.084] \\
gain scale, MLP norm output & 54 & 0.937 & 0.862 & 0.899 & $-$0.026 [$-$0.060, +0.001] \\
scale up, input pixels & 54 & 0.740 & 0.476 & 0.587 & $-$0.223 [$-$0.288, $-$0.162] \\
dropout, after both residual adds & 54 & 0.885 & 0.736 & 0.798 & $-$0.078 [$-$0.135, $-$0.026] \\
\bottomrule
\end{tabular}
\end{table*}


Masking whole tokens is the best perturbation at that site. It reaches \StaircaseOperatorsBeforeAttentionNormTokenMaskAuroc{} and flags \StaircaseOperatorsBeforeAttentionNormTokenMaskTprAtOneZeroPercent{} of the triggered inputs at a 10\% budget, against \StaircaseOperatorsBeforeAttentionNormChannelMaskAuroc{} for masking whole channels, \StaircaseOperatorsBeforeAttentionNormAuroc{} for dropout and \StaircaseOperatorsBeforeAttentionNormGaussianAuroc{} for Gaussian noise. This is PSBD-TM, and it reads below chance on \StaircaseOperatorsBeforeAttentionNormTokenMaskBelowChance{} of \StaircaseOperatorsBeforeAttentionNormTokenMaskN{} models against \StaircaseOperatorsPostResidualBelowChance{} for PSBD-RD.

Token masking on the attention branch output, just before its residual addition, is its twin. It reads \StaircaseOperatorsBeforeAttentionResidualTokenMaskAuroc{}, a paired difference of \StaircaseOperatorsBeforeAttentionResidualTokenMaskGain{} [\StaircaseOperatorsBeforeAttentionResidualTokenMaskGainLow{}, \StaircaseOperatorsBeforeAttentionResidualTokenMaskGainHigh{}] that we cannot separate from zero, so the two sites on either side of the attention block behave alike on ViT. We recommend the attention input for three reasons. It is the site whose mechanism \cref{sec:mechanism} explains, it is the more stable of the two across training seeds, and on Swin-S it leads its twin by \SwinGainRecommendedMinusTwin{} [\SwinGainRecommendedMinusTwinLow{}, \SwinGainRecommendedMinusTwinHigh{}] (\cref{sec:swin}). Token masking on the stream after the attention addition, by contrast, falls to \StaircaseOperatorsAfterAttentionResidualTokenMaskAuroc{}. The three sites differ only in whether the residual stream itself is masked, and \cref{sec:mechanism} explains why that matters.

What the perturbation is worth depends on which side of a LayerNorm it lands on rather than on the site. Injected before a LayerNorm, Gaussian noise loses to token masking at both sites, by \StaircaseOperatorsBeforeAttentionNormGaussianGain{} at the attention input and by \StaircaseOperatorsBeforeMlpNormGaussianGain{} at the MLP input. Injected after the MLP's LayerNorm the same noise reads \StaircaseOperatorsBeforeMlpGaussianAuroc{}, above token masking before it, so what matters is the side of the norm and not the site. \Cref{sec:mechanism:layernorm} explains both this and why whole tokens beat coordinates.

\subsection{Across attacks, poison rates and datasets}
\label{sec:results:headline}

\Cref{tab:results-rate-1,tab:results-rate-5,tab:results-rate-10} in \cref{app:results} give PSBD-TM and PSBD-RD on every model at \PanelRates{} poisoning, with ASR, clean accuracy, AUROC and TPR@FPR at both budgets. \Cref{tab:gains} summarizes the paired gain by poison rate and dataset.

% generated by scripts/paper/tab_gains.py from results/coverage/coverage.json (56 successful cells), results/*/psbd_metrics.json at 2026-09-29T18:07:49.405572+00:00, commit c6325b925e68b12d75a1f415faa829e94c9b986b-dirty. Do not edit by hand.
\begin{table}[htbp]
\centering
\small
\caption{Paired AUROC gain of token masking at the attention input over dropout after both residual adds, at the adaptive rule and the headline quantile, with 95\% bootstrap intervals over models from 5000 resamples. Every row is a paired difference on the models carrying both placements.}
\label{tab:gains}
\begin{adjustbox}{max width=\linewidth}
\begin{tabular}{lrrl}
\toprule
subset & models & gain & 95\% interval \\
\midrule
all models & 54 & +0.078 & [+0.026, +0.135] \\
\multicolumn{4}{l}{\emph{by poison rate}} \\
1\% poisoning & 17 & +0.150 & [+0.044, +0.279] \\
5\% poisoning & 19 & +0.077 & [+0.018, +0.149] \\
10\% poisoning & 18 & +0.012 & [$-$0.065, +0.097] \\
\multicolumn{4}{l}{\emph{by dataset}} \\
CIFAR-10 & 15 & +0.088 & [$-$0.048, +0.231] \\
Tiny ImageNet & 14 & +0.012 & [$-$0.007, +0.039] \\
GTSRB & 13 & +0.129 & [+0.020, +0.272] \\
CIFAR-100 & 12 & +0.088 & [+0.012, +0.183] \\
\bottomrule
\end{tabular}
\end{adjustbox}
\end{table}


The gain grows as the poison rate falls, from \GainsHighestRate{} at \PanelRateHigh{} to \GainsLowestRate{} at \PanelRateLow{}. This is the regime where a defender most needs help, since an attacker who can poison only \PanelRateLow{} of the data plants a weaker shortcut. The gain is positive on \GainsPositiveDatasets{} of \GainsDatasets{} datasets, large on CIFAR-100 and GTSRB and small on Tiny ImageNet, where PSBD-RD is already strong on BadNets, Blend and LF. Removing any single attack leaves a gain of at least \GainsLeaveOneAttackOut{}, with a lower bootstrap bound of \GainsLeaveOneAttackOutLow{}, so no single attack carries the result. Per attack, PSBD-TM stays at or above \GainsStrongAuroc{} on every dataset and rate only for \GainsStrongAttacks{}. Every other attack has at least one model below that, and the floor is \HeadlineFloorAuroc{} on the two models \cref{sec:limitations} discusses.

\Cref{app:distributions} shows the PSU distributions behind these numbers on six models.
